\chapter{Theory}
\label{ch:theory}

\section{Finite-strain hyperelasticity}

Let \(\bm X\) denote a material point in the reference configuration and
\(\bm x=\bm\varphi(\bm X)\) its current position. The deformation gradient and
local volume ratio are
\begin{equation}
  \F=\frac{\partial\bm\varphi}{\partial\bm X},
  \qquad J=\det\F>0.
\end{equation}
For a hyperelastic material, the first-Piola stress and total-Lagrange
material tangent follow from a stored energy,
\begin{equation}
  P_{iJ}=\frac{\partial\Psi}{\partial F_{iJ}},
  \qquad
  A_{iJkL}
  =\frac{\partial^2\Psi}{\partial F_{iJ}\partial F_{kL}}.
  \label{eq:stress-tangent}
\end{equation}
Because \(\F\) and \(\Pstress\) are generally unsymmetric, the implementation
retains all nine components and uses
\begin{equation}
  A^{9\times9}_{3i+J,\,3k+L}=A_{iJkL}.
  \label{eq:flattening}
\end{equation}
This is not a \(6\times6\) Voigt representation of symmetric stress and
strain measures.

In a total-Lagrangian finite-element formulation, equilibrium follows from
stationarity of the potential with respect to an admissible displacement
variation \(\delta\bm u\):
\begin{equation}
  \int_{\Omega_0} \Pstress(\F):\nabla_X\delta\bm u\,\dd V
  = \delta W_{\mathrm{ext}}.
  \label{eq:weak-equilibrium}
\end{equation}
The internal contribution to the Newton matrix is the derivative of this
residual with respect to displacement coefficients. It contains
\(\Aten=\partial\Pstress/\partial\F\) contracted with the trial and test
gradients. Returning an energy value alone is therefore insufficient for a
nonlinear FE material interface: its first and second derivatives must use
the same quadrature locations and tensor ordering.

The Phase~1 analytical reference is the compressible neo-Hookean energy
\begin{equation}
  \Psi_{\mathrm{NH}}(\F)
  =\frac{\mu}{2}(I_1-3)-\mu\ln J
   +\frac{\lambda}{2}(\ln J)^2,
  \qquad I_1=\operatorname{tr}(\F^\mathsf T\F),
  \label{eq:neo-hooke}
\end{equation}
following the finite-strain framework summarized by
Holzapfel~\cite{holzapfel2000}.
Its analytical first-Piola stress is
\begin{equation}
  \Pstress_{\mathrm{NH}}
  =\mu(\F-\F^{-\mathsf T})
   +\lambda\ln J\,\F^{-\mathsf T}.
  \label{eq:neo-hooke-stress}
\end{equation}
This closed-form response is used as the Phase~1 target and as an independent
reference for testing PyTorch differentiation. The purpose of that stage is
not to identify the composite material; it verifies the software path from a
scalar learned energy through stress and tangent evaluation to a converged
NGSolve equilibrium problem.

\section{Polyconvex neural energy}

Convexity directly in \(\F\) is too restrictive for finite elasticity.
Polyconvexity instead represents the energy as a convex function of suitable
minors of \(\F\) and provides an important existence framework
\cite{ball1977}. The neural model uses the isotropic feature vector
\begin{equation}
  \bm y(\F)=\left[I_1,I_2,J,-\ln J\right],
\end{equation}
where
\begin{equation}
  I_2=\frac12\left[
    I_1^2-\operatorname{tr}\!\left((\F^\mathsf T\F)^2\right)
  \right].
\end{equation}
In three dimensions, \(I_2=\|\operatorname{cof}\F\|^2\). The feature map thus
uses the squared norms of \(\F\) and its cofactor, the determinant, and a
volumetric term. The first two features describe distortional deformation;
\(J\) and \(-\ln J\) allow the network to represent the asymmetric energetic
cost of volume change. The admissibility condition \(J>0\) is essential
because both the logarithm and the elastic model require an
orientation-preserving deformation.

An ICNN architecture~\cite{amos2017} constrains relevant weights through
\(W^+=\operatorname{softplus}(\theta)>0\). In the multiscale surrogate,
\(\vfrac\) acts as a non-convex context variable while convexity is retained
with respect to the deformation features. This construction follows the
physics-constrained neural hyperelasticity literature
\cite{kalina2022,linden2023}; it does not by itself establish positive
definiteness of the tangent at every deformation.

For each fixed context, nonnegative hidden-to-hidden and feature-to-hidden
weights together with the convex, nondecreasing softplus activation preserve
convexity of the scalar network in its feature inputs. In Phase~2, volume
fraction changes the learned effective response as context; it is not a
spatial field solved by the macroscopic finite-element model. The architecture
therefore imposes a useful energy structure while numerical tests are still
needed for stability, admissibility, and accuracy over the intended
deformation range.

The raw network is normalized at the reference state:
\begin{equation}
  \widehat\Psi(\F,\vfrac)
  =
  \Psi_{\mathrm{raw}}(\F,\vfrac)
  -\Psi_{\mathrm{raw}}(\Iten,\vfrac)
  -p_I(\vfrac)\bigl(\det\F-1\bigr),
  \qquad
  p_I(\vfrac)=\frac{1}{3}
  \operatorname{tr}\Pstress_{\mathrm{raw}}(\Iten,\vfrac).
  \label{eq:normalization}
\end{equation}
For an isotropic raw model,
\(\Pstress_{\mathrm{raw}}(\Iten,\vfrac)=p_I(\vfrac)\Iten\), so the correction
gives \(\widehat\Psi(\Iten,\vfrac)=0\) and
\(\partial\widehat\Psi/\partial\F|_{\Iten}=\bm 0\) while remaining objective.
It is affine in the determinant feature and therefore does not change
convexity with respect to the independent polyconvex features. Since
\(\det\F\) is nonlinear in \(\F\), however, this correction does change the
Hessian with respect to \(\F\); the constitutive tangent is differentiated
from the normalized energy itself. Positive affine feature and context
scalings are estimated from training data only and are stored with the model;
inference uses the same scalings.

\section{Mixed near-incompressible matrix}

The PDMS matrix is represented by an isochoric neo-Hookean energy with an
independent pressure:
\begin{equation}
  \Pi(\F,p)
  =
  \frac{\mu}{2}\left(J^{-2/3}I_1-3\right)
  +p(J-1)-\frac{p^2}{2\kappa}.
  \label{eq:mixed-energy}
\end{equation}
Stationarity with respect to \(p\) gives
\(p=\kappa(J-1)\). A Taylor--Hood \(P_2/P_1\)
displacement--pressure pair is used to reduce the risk of volumetric locking
as \(\kappa/\mu\) becomes large. Locking also depends on mesh resolution,
quadrature, and the chosen finite-element approximation; the mixed affine
patch test verifies consistency but is not itself a mesh-convergence study.

Eliminating \(p\) analytically recovers the volumetric penalty
\(\frac{\kappa}{2}(J-1)^2\). Keeping pressure as an independent FE field
instead gives a mixed saddle-point system, so a low-order displacement field
does not have to represent the near-incompressibility constraint by itself.
The implemented displacement stress consists of the isochoric contribution
plus \(p\,\operatorname{cof}\F\), while the pressure equation enforces
\(J-1-p/\kappa=0\). The selected ratio \(\kappa/\mu=10^4\) represents a
nearly incompressible matrix; it does not remove the need to check element
order, mesh density, integration, and conditioning.

\section{Periodic homogenization and rigid particles}

The microscopic deformation is decomposed as
\begin{equation}
  \F(\bm X)=\Fbar+\nabla_{\!X}\bm w(\bm X),
\end{equation}
with periodic fluctuation \(\bm w\). Opposing boundary traces coincide, and
their tractions are anti-periodic in the weak equilibrium sense. The rigid
particle motion is
\begin{equation}
  \bm u_{\mathrm p}(\bm X)
  =
  \bm t+(\Rten-\Iten)(\bm X-\bm X_{\mathrm c}),
  \qquad
  \Rten\in\mathrm{SO}(3).
  \label{eq:rigid-motion}
\end{equation}
The total displacement is
\(\bm u(\bm X)=(\Fbar-\Iten)\bm X+\bm w(\bm X)\). On the cavity surface,
\(\bm w\) is constrained so that the particle undergoes a rigid motion while
the matrix follows the imposed cell deformation. The inclusion is not
approximated as an extremely stiff elastic phase: its surface has no
deformational degrees of freedom, and its free orientation is determined by
zero resultant moment. A constant translation shared by the periodic
fluctuation and particle is a gauge mode; the macroscopic deformation maps
the particle center.

The Hill--Mandel condition requires equality of macroscopic and volume-averaged
microscopic power:
\begin{equation}
  \Pbar:\dot{\Fbar}
  =
  \frac{1}{|\Omega|}
  \int_{\Omega_{\mathrm m}}
  \Pstress:\dot{\F}\,\dd V.
  \label{eq:hill-mandel}
\end{equation}
For the cavity representation of a rigid inclusion, work-conjugate
macrostress requires both matrix-volume stress and assembled particle-reaction
work. Raw constitutive boundary tractions are retained only as a diagnostic.
Periodic fluctuation contributions cancel on opposing faces. The remaining
rigid-particle constraint contributes through its assembled reaction. Using
only the stress integrated over the deformable matrix would omit this work
and need not produce a stress conjugate to \(\Fbar\).

\section{Condensed macroscopic tangent}

Let \(\bm z\) collect free fluctuation, pressure, and rotational coordinates.
The exact consistent macroscopic tangent is the Schur complement of the
discrete energy Hessian,
\begin{equation}
  \Abar
  =
  \frac{1}{|\Omega|}
  \left(
    \bm H_{FF}
    -\bm H_{Fz}\bm H_{zz}^{-1}\bm H_{zF}
  \right).
  \label{eq:schur}
\end{equation}
Here \(\bm H\) is the Hessian of the discrete mixed energy with respect to the
nine macroscopic coordinates and internal variables \(\bm z\). The formula
is the derivative of the energy after the internal state has relaxed to
equilibrium: each macroscopic perturbation induces a microscopic
displacement, pressure, and, for a freely rotating particle, rotation
increment. The resulting matrix is the Hessian of the discrete
homogenized energy divided by total cell volume. It should therefore be
major-symmetric up to linear-solver and floating-point tolerances.
NGSolve symbolically differentiates the discrete energy. Consequently,
\(\Abar\) is not constructed from finite perturbations of separately solved
RVE states.
